\section{Diseño del prototipo}

El prototipo de PuellaGame se diseñó utilizando una estructura modular por capas, con el objetivo de separar las diferentes responsabilidades del sistema y facilitar su desarrollo, mantenimiento y prueba. El sistema administra tres entidades principales: Sucursal, Premio y Cliente. Estas tres entidades utilizan un identificador de tipo entero y se encuentran implementadas bajo un contrato común mediante la interfaz \texttt{Identificable}. Esto permite que las entidades compartan la operación para obtener y establecer su identificador y facilita que otros componentes del sistema puedan trabajar con ellas de manera uniforme.

Para la persistencia de la información se decidió utilizar archivos CSV, manteniendo un archivo independiente para cada entidad. De esta manera, la información de las sucursales se almacena en \texttt{sucursales.csv}, la información de los premios en \texttt{premios.csv} y la información de los clientes en \texttt{clientes.csv}. Esta organización permite mantener separada la información correspondiente a cada entidad y facilita su lectura y almacenamiento.

La persistencia fue separada del resto del sistema mediante la capa \texttt{Repository}. Esta capa se encarga de realizar las operaciones sobre los archivos CSV, evitando que la interfaz o la lógica de negocio tengan que conocer directamente cómo se almacenan los datos. Para representar la estructura de cada archivo se utiliza \texttt{CsvSchema}, que define las columnas, la forma de obtener la llave de una entidad y las funciones necesarias para convertir una entidad a una fila CSV y reconstruir una entidad a partir de los datos almacenados.

Sobre la capa \texttt{Repository} se encuentra la capa \texttt{Service}, que concentra las operaciones de negocio del sistema. Esta capa permite realizar las operaciones de alta, consulta, modificación y eliminación de las tres entidades, además de la búsqueda mediante el identificador. De esta manera, la lógica relacionada con los casos de uso queda separada tanto del almacenamiento como de la interacción con el usuario.

La interacción con el usuario se realiza mediante una interfaz de consola. La capa \texttt{UI} contiene los menús y formularios necesarios para capturar la información de sucursales, premios y clientes. Se cuenta con un menú principal y con submenús CRUD para cada una de las entidades. Los formularios específicos permiten solicitar los datos correspondientes a cada entidad, mientras que componentes como \texttt{ConsoleInput} centralizan la lectura y validación de las entradas.

La validación de los datos se encuentra separada en una capa propia. Esta capa comprueba que los valores ingresados cumplan con las condiciones establecidas para el sistema, como utilizar enteros positivos para los identificadores, enteros no negativos para los puntos y existencias, textos no vacíos, correos electrónicos válidos y teléfonos de diez dígitos. Cuando una entrada no cumple con las condiciones esperadas, se utiliza \texttt{ValidationException} para representar el error.

De manera similar, se utiliza una capa de excepciones para comunicar diferentes tipos de errores. Por ejemplo, \texttt{EntityNotFoundException} representa el caso en el que una entidad solicitada no existe. Estas excepciones son manejadas por la interfaz de consola para mostrar al usuario mensajes de error comprensibles sin interrumpir la ejecución normal de la aplicación.

Finalmente, la clase \texttt{Main} se encarga de integrar los diferentes componentes del sistema. En ella se crean los repositorios para cada archivo CSV, los servicios correspondientes y los componentes de la interfaz de consola. De esta manera, las diferentes capas se conectan para formar la aplicación completa.

La arquitectura modular adoptada permite separar responsabilidades entre los distintos componentes del sistema. El modelo representa los datos, \texttt{Repository} administra la persistencia, \texttt{Service} concentra las operaciones de negocio, \texttt{UI} maneja la interacción con el usuario, \texttt{Validation} verifica las entradas y \texttt{Exception} permite representar y comunicar errores. Esta separación facilita que cada componente pueda desarrollarse y probarse de manera independiente, además de hacer más sencillo modificar una parte del sistema sin afectar directamente a las demás.
